\section*{Chapter \ref{ch:g10:earth-energy-budget} --- The Earth's Energy Budget}

\begin{solution}{exo:g10:earth-energy-budget:1}
The \omterm{def:g10:earth-energy-budget:insolation}{insolation} of a place is the power of sunlight received by a horizontal
square metre at the top of the atmosphere above it (averaged over a day or a
year). The \omterm{def:g10:earth-energy-budget:albedo}{albedo} of a surface is the fraction of the sunlight falling on it
that it reflects. Fresh snow (0.8 to 0.9) reflects far more than the open
ocean (0.06).
\end{solution}

\begin{solution}{exo:g10:earth-energy-budget:2}
Reflected: $100/340 \approx \qty{29}{\percent}$; absorbed by the air:
$80/340 \approx \qty{24}{\percent}$; absorbed by the surface: $160/340
\approx \qty{47}{\percent}$.
\end{solution}

\begin{solution}{exo:g10:earth-energy-budget:3}
Asphalt: $800 \times (1 - 0.04) = \qty{768}{W}$. Snow: $800 \times (1 -
0.85) = \qty{120}{W}$, more than six times less.
\end{solution}

\begin{solution}{exo:g10:earth-energy-budget:4}
Evaporation: \qty{82}{W/m^2}, against \qty{21}{W/m^2} for the warming of the
air. The water that evaporates takes its \omterm{def:g10:earth-energy-budget:latent-heat}{latent heat} from the sea; the vapour
rises with the air, and where it condenses into cloud droplets, often
kilometres higher, it gives that energy back to the air around it.
\end{solution}

\begin{solution}{exo:g10:earth-energy-budget:5}
In December the Northern Hemisphere leans away from the Sun: at the pole the
Sun stays below the horizon all day, so no sunlight arrives. On the December
solstice this is so everywhere poleward of $90^\circ - 23.44^\circ =
\qty{66.56}{\degree}$ N.
\end{solution}

\begin{solution}{exo:g10:earth-energy-budget:6}
$1361 \sin 90^\circ = \qty{1361}{W}$; $1361 \sin 30^\circ =
\qty{680.5}{W}$; $1361 \sin 10^\circ \approx \qty{236}{W}$. Half of $S_0$
needs $\sin h = 0.5$, that is $h = \qty{30}{\degree}$.
\end{solution}

\begin{solution}{exo:g10:earth-energy-budget:7}
Reflected 100, absorbed by the air 80, absorbed by the surface 160:
$100 + 80 + 160 = 340$, all the sunlight is accounted for. At the surface,
$185 = 25 + 160$: what is not reflected is absorbed.
\end{solution}

\begin{solution}{exo:g10:earth-energy-budget:8}
Disc: $\pi \times (6.371\times10^6)^2 \approx \qty{1.275e14}{m^2}$; power:
$1361 \times 1.275\times10^{14} \approx \qty{1.74e17}{W}$. Sphere: $4\pi
R^2 \approx \qty{5.10e14}{m^2}$. Average: $1.74\times10^{17}/5.10\times
10^{14} \approx \qty{340}{W/m^2}$, a quarter of $S_0$.
\end{solution}

\begin{solution}{exo:g10:earth-energy-budget:9}
$0.7 \times 5.10\times10^{14} \times 3.16\times10^7 \approx
\qty{1.1e22}{J}$ in a year.
\end{solution}

\begin{solution}{exo:g10:earth-energy-budget:10}
The clouds reflect $100 - 53 = \qty{47}{W/m^2}$ more sunlight, which is
therefore not absorbed. They hold back $267 - 239 = \qty{28}{W/m^2}$ of
infrared. On balance the Earth loses $47 - 28 = \qty{19}{W/m^2}$: the clouds,
as a whole, cool it.
\end{solution}

\begin{solution}{exo:g10:earth-energy-budget:11}
From about 10 May to the end of July, that is in May, June and July. On the
equator the Sun passes overhead at noon at the two equinoxes, in March and
September; at the solstices it passes \qty{23.44}{\degree} away from the
zenith, so the equator receives a little less in June and in December.
\end{solution}

\begin{solution}{exo:g10:earth-energy-budget:12}
Each second: $82/(2.47\times10^6) \approx \qty{3.32e-5}{kg}$. Each day:
$3.32\times10^{-5} \times 86\,400 \approx \qty{2.9}{kg}$, a layer of about
\qty{2.9}{mm}. Each year: $3.32\times10^{-5} \times 3.16\times10^7 \approx
\qty{1050}{kg}$, about \qty{1.05}{m} of water.
\end{solution}

\begin{solution}{exo:g10:earth-energy-budget:13}
$S_0 \sin\delta = S_0/\pi$ gives $\sin\delta = 1/\pi \approx 0.318$, so
$\delta \approx \qty{18.6}{\degree}$. Then $\sin\lambda = 0.318/\sin 23.44^\circ
\approx 0.318/0.398 \approx 0.800$: $\lambda \approx \qty{53.15}{\degree}$ or
$180^\circ - 53.15^\circ = \qty{126.85}{\degree}$. Between the two, $\lambda$
grows by \qty{73.7}{\degree}, that is $73.7/360$ of a year, about
\num{75} days, from mid-May to late July.
\end{solution}

\begin{solution}{exo:g10:earth-energy-budget:14}
The ice absorbs $1 - 0.6 = 0.40$ of the sunlight, the ocean $1 - 0.06 =
0.94$. Extra power: $(0.94 - 0.40) \times 300 \times 10^{12} \approx
\qty{1.6e14}{W}$, more than all the photosynthesis of the planet. The extra
energy warms the water, which melts more ice at the edges of the floes,
which exposes more dark water: the change feeds itself.
\end{solution}

\begin{solution}{exo:g10:earth-energy-budget:15}
Closest: $1361 \times (149.6/147.1)^2 \approx \qty{1408}{W/m^2}$; farthest:
$1361 \times (149.6/152.1)^2 \approx \qty{1317}{W/m^2}$; about
\qty{7}{\percent} apart. At \qty{60}{\degree} N the \omterm{def:g10:earth-energy-budget:insolation}{insolation} is twenty
times larger in June than in December: a change of a few per cent cannot
cause it. Besides, the Earth is closest to the Sun in early January, during
the northern winter, and the two hemispheres, at the same distance from the
Sun, have opposite seasons: the tilt of the axis, not the distance, makes
the seasons.
\end{solution}

\begin{solution}{pb:g10:earth-energy-budget:1}
\textbf{1.} $\pi \times (6.371\times10^6)^2 \approx \qty{1.275e14}{m^2}$.

\textbf{2.} $1361 \times 1.275\times10^{14} \approx \qty{1.735e17}{W}$.

\textbf{3.} $4\pi R^2 \approx \qty{5.101e14}{m^2}$.

\textbf{4.} $1.735\times10^{17}/5.101\times10^{14} \approx \qty{340}{W/m^2}$,
that is $S_0/4$.

\textbf{5.} Equator: \qty{1361}{W}; \qty{60}{\degree} N: $1361 \sin
30^\circ = \qty{680.5}{W}$.

\textbf{6.} $100/340 \approx 0.294$.

\textbf{7.} $340 - 100 = \qty{240}{W/m^2}$; over the Earth, $240 \times
5.101\times10^{14} \approx \qty{1.224e17}{W}$.

\textbf{8.} $53/340 \approx 0.156$. Clouds raise the \omterm{def:g10:earth-energy-budget:albedo}{albedo} by a factor
$100/53 \approx 1.9$: they almost double it.

\textbf{9.} $(240/5.67\times10^{-8})^{1/4} \approx \qty{255}{K}$, about
\qty{33}{K} below the surface's \qty{288}{K}.

\textbf{10.} $160 + 342 = \qty{502}{W/m^2}$.

\textbf{11.} $82 + 21 + 398 = \qty{501}{W/m^2}$. The surface gains about
\qty{1}{W/m^2}: the imbalance of \qty{0.7}{W/m^2}, rounded.

\textbf{12.} $342/160 \approx 2.1$. The air sends the surface more than
twice as much energy as infrared as the sunlight it absorbs: the greenhouse
effect.

\textbf{13.} $82 \times 3.156\times10^7/(2.47\times10^6) \approx
\qty{1050}{kg}$, a layer about \qty{1.05}{m} thick.

\textbf{14.} $0.7 \times 5.101\times10^{14} \approx \qty{3.6e14}{W}$.

\textbf{15.} $2805/72.07 \approx \qty{38.9}{kJ}$ per gram of carbon.

\textbf{16.} $104.9\times10^9$ tonnes is $1.049\times10^{17}$~g;
$1.049\times10^{17} \times 38.9\times10^3 \approx \qty{4.08e21}{J}$.

\textbf{17.} $4.08\times10^{21}/3.156\times10^7 \approx \qty{1.29e14}{W}$.

\textbf{18.} $1.29\times10^{14}/5.101\times10^{14} \approx
\qty{0.25}{W/m^2}$.

\textbf{19.} $0.7/0.25 \approx 2.8$: the Earth's imbalance is almost three
times the power that photosynthesis stores.

\textbf{20.} $1.29\times10^{14}/1.224\times10^{17} \approx 0.0011$:
photosynthesis stores \textbf{about \qty{0.1}{\percent}} of the sunlight
the Earth absorbs --- one thousandth.
\end{solution}
